% Generated from lessons/0009-2-upper-lower-limits.html; edit the HTML source, then rerun the converter.
\section{上极限与下极限}
\lessonmark{上极限与下极限}

这一节讨论一个很重要的考试问题：当数列不收敛时，我们怎样精确描述它“最后到底在什么范围里晃”。核心工具就是极限点、上极限 \(\overline{\lim}\) 和下极限 \(\underline{\lim}\)。这一节既是数列章节的深化，也是后面研究级数、函数列、积分判别时常用的语言工具。

\subsection*{本节主线}

\begin{conceptbox}
\textbf{第一步：看子列}
不收敛的数列，常常不同子列有不同极限。先收集所有可能的子列极限。
\end{conceptbox}

\begin{conceptbox}
\textbf{第二步：取最大最小}
所有极限点的最大值叫上极限，最小值叫下极限。
\end{conceptbox}

\begin{conceptbox}
\textbf{第三步：转成判题模板}
上极限等于“最终不超过某值，并且无穷多次逼近这个值”。这是最实用的考试版定义。
\end{conceptbox}

\begin{notebox}
一句话记忆：上极限看“最后还能反复达到多高”，下极限看“最后还能反复掉到多低”。
\end{notebox}

\subsection*{一、为什么需要上极限与下极限}

如果数列 \(\{x_n\}\) 收敛，我们只需要一个数 \(\lim x_n\) 来描述它。但很多数列并不收敛，例如

\[
      x_n=(-1)^n,\qquad
      x_n=\sin\frac{n\pi}{3},\qquad
      x_n=n+(-1)^n.
    \]

这些数列虽然不收敛，但也不是“完全没规律”：它们往往会沿某些子列逼近某些特定的值。于是我们先研究：

\[
      \text{这个数列有哪些子列极限？}
    \]

再从这些子列极限里挑出最大的和最小的，这就得到上极限和下极限。

\subsection*{二、极限点：先收集所有子列极限}

\subsubsection*{定义 9.2.1：极限点（子列极限）}

设 \(\{x_n\}\) 是有界数列。若存在它的一个子列 \(\{x_{n_k}\}\)，使得

\[
      \lim_{k\to\infty}x_{n_k}=\xi,
    \]

则称 \(\xi\) 是数列 \(\{x_n\}\) 的一个\textbf{极限点}，也叫\textbf{子列极限}。

\subsubsection*{等价直观}

\(\xi\) 是极限点，等价于：对任意 \(\varepsilon>0\)，区间 \((\xi-\varepsilon,\xi+\varepsilon)\) 里总能找到无穷多个数列项。

\begin{notebox}
这里要特别分清：
\end{notebox}

\begin{itemize}
 \item \textbf{“最终都进去”} 是整个数列收敛到 \(\xi\)；
 \item \textbf{“总有无穷多个进去”} 只是说 \(\xi\) 是一个极限点。
 \end{itemize}

\subsubsection*{例子}

\begin{conceptbox}
\textbf{\(x_n=(-1)^n\)}
偶数子列恒为 \(1\)，奇数子列恒为 \(-1\)。所以极限点是 \(\{-1,1\}\)。
\end{conceptbox}

\begin{conceptbox}
\textbf{\(x_n=\frac{(-1)^n}{n}\)}
虽然正负交替，但所有子列都趋于 \(0\)，所以唯一极限点是 \(0\)。
\end{conceptbox}

\subsection*{三、上极限与下极限的定义}

记数列 \(\{x_n\}\) 的全体极限点组成的集合为

\[
      E=\{\xi:\ \xi \text{ 是 } \{x_n\}\text{ 的极限点}\}.
    \]

若 \(\{x_n\}\) 有界，则由 Bolzano-Weierstrass 定理，\(E\neq\varnothing\)，而且 \(E\) 也是有界集，所以

\[
      H=\sup E,\qquad h=\inf E
    \]

都存在。

\subsubsection*{定理 9.2.1：上确界和下确界本身也是极限点}

\[
      H\in E,\qquad h\in E.
    \]

因此实际上

\[
      H=\max E,\qquad h=\min E.
    \]

\paragraph*{证明思路}

只说 \(H\)。因为 \(H=\sup E\)，可取 \(\xi_k\in E\)，使

\[
      H-\frac1k<\xi_k\le H.
    \]

而每个 \(\xi_k\) 本身是极限点，所以都能在原数列里找到很多项非常靠近 \(\xi_k\)。从中递增地抽出一项 \(x_{n_k}\)，使

\[
      |x_{n_k}-\xi_k|<\frac1k.
    \]

于是

\[
      |x_{n_k}-H|
      \le
      |x_{n_k}-\xi_k|+|\xi_k-H|
      <
      \frac1k+\frac1k
      \to0.
    \]

故 \(x_{n_k}\to H\)，所以 \(H\) 也是极限点。\(h\) 同理。

\subsubsection*{定义 9.2.2：上极限与下极限}

\[
      \overline{\lim}_{n\to\infty}x_n:=H=\max E,
      \qquad
      \underline{\lim}_{n\to\infty}x_n:=h=\min E.
    \]

\begin{notebox}
所以：上极限不是随便取个大数，它必须真的是某个子列的极限；下极限也是一样。
\end{notebox}

\subsection*{四、它和普通极限的关系}

\subsubsection*{定理 9.2.2：收敛判别}

对有界数列 \(\{x_n\}\)，

\[
      \{x_n\}\text{ 收敛}
      \iff
      \overline{\lim}_{n\to\infty}x_n
      =
      \underline{\lim}_{n\to\infty}x_n.
    \]

并且在这种情形下，三者都等于同一个极限：

\[
      \overline{\lim}_{n\to\infty}x_n
      =
      \underline{\lim}_{n\to\infty}x_n
      =
      \lim_{n\to\infty}x_n.
    \]

\paragraph*{为什么}

\begin{itemize}
 \item 如果 \(x_n\to a\)，那么任何子列也都趋于 \(a\)，所以唯一极限点就是 \(a\)。
 \item 如果上极限和下极限相等，说明全部子列极限只有一个值，数列就只能收敛到这个值。
 \end{itemize}

\subsubsection*{无界情形的扩充}

对一般实数列，也允许 \(+\infty\) 和 \(-\infty\) 作为“广义极限点”。于是上极限、下极限可以取扩充实数值，且仍有：

\[
      \lim_{n\to\infty}x_n \text{ 存在（有限或 } \pm\infty\text{）}
      \iff
      \overline{\lim}_{n\to\infty}x_n
      =
      \underline{\lim}_{n\to\infty}x_n.
    \]

\subsection*{五、最实用的考试判别法}

\subsubsection*{定理 9.2.3：上极限的 \(\varepsilon\)-刻画}

对有界数列 \(\{x_n\}\)，

\[
      \overline{\lim}_{n\to\infty}x_n=H
    \]

当且仅当对任意 \(\varepsilon>0\)，都有：

\begin{enumerate}
 \item 存在 \(N\)，当 \(n>N\) 时，\(x_n<H+\varepsilon\)；
 \item 有无穷多个 \(n\) 满足 \(x_n>H-\varepsilon\)。
 \end{enumerate}

同理，

\[
      \underline{\lim}_{n\to\infty}x_n=h
    \]

当且仅当对任意 \(\varepsilon>0\)，都有：

\begin{enumerate}
 \item 存在 \(N\)，当 \(n>N\) 时，\(x_n>h-\varepsilon\)；
 \item 有无穷多个 \(n\) 满足 \(x_n<h+\varepsilon\)。
 \end{enumerate}

\begin{notebox}
\textbf{这一定理一定要会翻成口语：}\\  上极限 \(H\) = “从某一项以后，数列再也不会明显高过 \(H\)；但它又会无穷多次逼近 \(H\) 的下方。”\\  下极限 \(h\) = “从某一项以后，数列再也不会明显低过 \(h\)；但它又会无穷多次逼近 \(h\) 的上方。”
\end{notebox}

\paragraph*{证明思路为什么成立}

如果 \(H\) 是最大的极限点，那么：

\begin{itemize}
 \item 不可能有一串尾项总是大于 \(H+\varepsilon\)，否则能抽出一个子列，得到比 \(H\) 更大的极限点；
 \item 也不可能只有有限多项大于 \(H-\varepsilon\)，否则 \(H\) 就不可能被子列逼近。
 \end{itemize}

\subsection*{六、运算性质：什么时候能像普通极限那样算}

\subsubsection*{1. 和的运算}

\[
      \overline{\lim}(x_n+y_n)
      \le
      \overline{\lim}x_n+\overline{\lim}y_n,
    \]

\[
      \underline{\lim}(x_n+y_n)
      \ge
      \underline{\lim}x_n+\underline{\lim}y_n.
    \]

一般只是\textbf{不等式}，不要直接写等号。

但如果 \(\lim x_n=x\) 存在，那么

\[
      \overline{\lim}(x_n+y_n)=x+\overline{\lim}y_n,
      \qquad
      \underline{\lim}(x_n+y_n)=x+\underline{\lim}y_n.
    \]

\subsubsection*{2. 积的运算}

若 \(x_n\ge0,\ y_n\ge0\)，则

\[
      \overline{\lim}(x_ny_n)
      \le
      \overline{\lim}x_n\cdot\overline{\lim}y_n,
    \]

\[
      \underline{\lim}(x_ny_n)
      \ge
      \underline{\lim}x_n\cdot\underline{\lim}y_n.
    \]

若 \(\lim x_n=x\) 且 \(0<x<+\infty\)，则

\[
      \overline{\lim}(x_ny_n)=x\cdot\overline{\lim}y_n,
      \qquad
      \underline{\lim}(x_ny_n)=x\cdot\underline{\lim}y_n.
    \]

\begin{notebox}
考试提醒：上极限、下极限的四则运算没有普通极限那么自由。只有“其中一个真的收敛了”时，公式才最稳。
\end{notebox}

\subsection*{七、最常用的计算工具：尾上确界与尾下确界}

对数列 \(\{x_n\}\)，定义尾部的上确界与下确界：

\[
      b_n=\sup\{x_n,x_{n+1},x_{n+2},\ldots\},
      \qquad
      a_n=\inf\{x_n,x_{n+1},x_{n+2},\ldots\}.
    \]

则 \(\{b_n\}\) 单调不增，\(\{a_n\}\) 单调不减，所以在扩充实数意义下都收敛。于是可定义

\[
      \limsup_{n\to\infty}x_n:=\lim_{n\to\infty}b_n,
      \qquad
      \liminf_{n\to\infty}x_n:=\lim_{n\to\infty}a_n.
    \]

\subsubsection*{定理 9.2.6：两种定义完全等价}

\[
      \limsup_{n\to\infty}x_n=\overline{\lim}_{n\to\infty}x_n,
      \qquad
      \liminf_{n\to\infty}x_n=\underline{\lim}_{n\to\infty}x_n.
    \]

也就是说：

\[
      \limsup_{n\to\infty}x_n
      =
      \inf_{n\ge1}\sup_{k\ge n}x_k,
      \qquad
      \liminf_{n\to\infty}x_n
      =
      \sup_{n\ge1}\inf_{k\ge n}x_k.
    \]

\begin{conceptbox}
\textbf{什么时候用极限点法}
数列明显能拆成几个子列，且每个子列的极限容易看出来。
\end{conceptbox}

\begin{conceptbox}
\textbf{什么时候用尾确界法}
题目要证明一般性质，或者数列结构不适合直接拆子列时。
\end{conceptbox}

\subsection*{八、例题串讲}

\subsubsection*{例 1：\(\displaystyle x_n=\cos\frac{2n\pi}{5}\)}

它只会取五个值，而且周期为 \(5\)。所以极限点就是这五个值本身：

\[
      \left\{\cos\frac{2\pi}{5},\cos\frac{4\pi}{5},\cos\frac{6\pi}{5},\cos\frac{8\pi}{5},\cos 2\pi\right\}.
    \]

其中最大的是 \(1\)，最小的是 \(-\cos\frac{\pi}{5}\)。因此

\[
      \overline{\lim}_{n\to\infty}x_n=1,
      \qquad
      \underline{\lim}_{n\to\infty}x_n=-\cos\frac{\pi}{5}.
    \]

\subsubsection*{例 2：\(\displaystyle x_n=n^{\frac{1-(-1)^n}{2}}\)}

把奇偶拆开：

\begin{itemize}
 \item 若 \(n=2k\)，则 \(\frac{1-(-1)^n}{2}=0\)，故 \(x_{2k}=1\)；
 \item 若 \(n=2k-1\)，则 \(\frac{1-(-1)^n}{2}=1\)，故 \(x_{2k-1}=2k-1\to+\infty\)。
 \end{itemize}

所以极限点是 \(1\) 和 \(+\infty\)，于是

\[
      \overline{\lim}_{n\to\infty}x_n=+\infty,
      \qquad
      \underline{\lim}_{n\to\infty}x_n=1.
    \]

\subsubsection*{例 3：\(\displaystyle x_n=\sqrt[n]{n+1}+\sin\frac{n\pi}{3}\)}

前一项 \(\sqrt[n]{n+1}\to1\)，后一项在

\[
      \left\{\frac{\sqrt3}{2},\frac{\sqrt3}{2},0,-\frac{\sqrt3}{2},-\frac{\sqrt3}{2},0\right\}
    \]

之间循环。因此

\[
      \overline{\lim}_{n\to\infty}x_n=1+\frac{\sqrt3}{2},
      \qquad
      \underline{\lim}_{n\to\infty}x_n=1-\frac{\sqrt3}{2}.
    \]

\begin{notebox}
\textbf{求上极限下极限的标准动作：}\\  1. 先看能不能按奇偶、模周期、分段形式拆子列；\\  2. 求每个子列极限；\\  3. 最大子列极限就是上极限，最小子列极限就是下极限；\\  4. 若有一支子列奔向 \(+\infty\)，则上极限就是 \(+\infty\)。
\end{notebox}

\subsection*{九、考试判断卡}

\begin{center}
\small
\begin{tabularx}{\linewidth}{|>{\raggedright\arraybackslash}p{0.18\linewidth}|>{\raggedright\arraybackslash}p{0.42\linewidth}|>{\raggedright\arraybackslash}X|}
\hline
\textbf{题型} & \textbf{先想什么} & \textbf{常用结论} \\ \hline
求 \(\overline{\lim},\underline{\lim}\) & 拆子列，找全部可能子列极限 & 上极限 = 最大极限点；下极限 = 最小极限点 \\ \hline
证明 \(\overline{\lim}x_n=H\) & 用“最终不超过 \(H+\varepsilon\)”和“无穷多项大于 \(H-\varepsilon\)” & 定理 9.2.3 \\ \hline
证明数列收敛 & 算出上极限和下极限 & 若二者相等，则数列收敛 \\ \hline
和、积里的 \(\limsup\) & 先问：是否有一个因子真的收敛？ & 若有，再用等式公式；否则通常只有不等式 \\ \hline
\end{tabularx}
\end{center}

\subsection*{十、即时判断练习}

\begin{quickcheck}
若 \(\overline{\lim}x_n=\underline{\lim}x_n\)，则 \(\{x_n\}\) 如何？

\tcblower
\textbf{答案：}
\end{quickcheck}

\begin{quickcheck}
\(\overline{\lim}x_n=H\) 最关键的两句话是？

\tcblower
\textbf{答案：}
\end{quickcheck}

\begin{quickcheck}
若 \(x_{2n}\to 3,\ x_{2n-1}\to -1\)，则 \(\overline{\lim}x_n\) 是？

\tcblower
\textbf{答案：}
\end{quickcheck}

\subsection*{十一、课后题训练}

下面按教材第九章第 2 节习题整理。题目完整保留，提示只给解题入口。写作业时建议先独立判断，再展开提示对照。

\begin{exercisebox}
\textbf{习题 1 \badge{求上极限与下极限}}

求下列数列的上极限与下极限：

\[
      \begin{aligned}
      &(1)\ x_n=\frac{n}{2n+1}\cos\frac{2n\pi}{5};\\
      &(2)\ x_n=n+(-1)^n\frac{n^2+1}{n};\\
      &(3)\ x_n=-n\bigl[(-1)^n+2\bigr];\\
      &(4)\ x_n=\sqrt[n]{n+1}+\sin\frac{n\pi}{3};\\
      &(5)\ x_n=2(-1)^{n+1}+3(-1)^{\frac{n(n-1)}2}.
      \end{aligned}
      \]

\begin{hintbox}
\begin{itemize}
 \item (1) \(\frac{n}{2n+1}\to \frac12\)，只需分析余下的周期因子 \(\cos\frac{2n\pi}{5}\)。
 \item (2) 先把 \(\frac{n^2+1}{n}=n+\frac1n\) 展开，再拆奇偶。
 \item (3) 奇偶两支都趋向 \(-\infty\)。
 \item (4) \(\sqrt[n]{n+1}\to1\)，后一项是周期数列。
 \item (5) 两个符号项都有周期，列出前四项就够了。
 \end{itemize}
\end{hintbox}
\end{exercisebox}

\begin{exercisebox}
\textbf{习题 2 \badge{符号与数乘}}

证明：

\[
      (1)\ \overline{\lim}_{n\to\infty}(-x_n)
      =
      -\,\underline{\lim}_{n\to\infty}x_n;
      \]

\[
      (2)\ \overline{\lim}_{n\to\infty}(c x_n)
      =
      \begin{cases}
      c\,\overline{\lim}_{n\to\infty}x_n, & c>0,\\[4pt]
      c\,\underline{\lim}_{n\to\infty}x_n, & c<0.
      \end{cases}
      \]

\begin{hintbox}
“乘上 \(-1\)” 会把所有极限点的大小顺序颠倒；“乘上正数”保持顺序，“乘上负数”反转顺序。
\end{hintbox}
\end{exercisebox}

\begin{exercisebox}
\textbf{习题 3 \badge{和的下极限}}

证明：

\[
      (1)\ \underline{\lim}_{n\to\infty}(x_n+y_n)
      \ge
      \underline{\lim}_{n\to\infty}x_n
      +
      \underline{\lim}_{n\to\infty}y_n;
      \]

\[
      (2)\ \text{若 } \lim_{n\to\infty}x_n \text{ 存在，则 }
      \underline{\lim}_{n\to\infty}(x_n+y_n)
      =
      \lim_{n\to\infty}x_n
      +
      \underline{\lim}_{n\to\infty}y_n.
      \]

\begin{hintbox}
直接对下极限用定理 9.2.3；或者先利用

\[
          \underline{\lim}(x_n+y_n)
          =
          -\,\overline{\lim}(-x_n-y_n)
        \]

再转成上极限公式。
\end{hintbox}
\end{exercisebox}

\begin{exercisebox}
\textbf{习题 4 \badge{负极限与乘法}}

证明：若 \(\lim_{n\to\infty}x_n=x\)，且 \(-\infty<x<0\)，则

\[
        \overline{\lim}_{n\to\infty}(x_ny_n)
        =
        x\,\underline{\lim}_{n\to\infty}y_n,
      \]

\[
        \underline{\lim}_{n\to\infty}(x_ny_n)
        =
        x\,\overline{\lim}_{n\to\infty}y_n.
      \]

\begin{hintbox}
负数乘法会把大小顺序翻转。可以先令 \(x_n=-u_n\)，其中 \(u_n\to -x>0\)，再套正数情形。
\end{hintbox}
\end{exercisebox}

\subsection*{十二、这一节的选题建议}

\begin{conceptbox}
\textbf{必做}
习题 1 的五小问全部要做；这是“拆子列找极限点”的标准训练。
\end{conceptbox}

\begin{conceptbox}
\textbf{证明题重点}
习题 2(1)、习题 3(1)。它们最能帮助你熟悉上极限下极限的运算规则。
\end{conceptbox}

\begin{conceptbox}
\textbf{压轴理解}
习题 4 是“负数会翻转上下极限”的典型题，很适合记进错题集或好题集。
\end{conceptbox}

来源参考：教材《数学分析》第三版下册第九章 §2（上极限与下极限）。这一节如果你卡在某个“为什么是这个上极限”，通常都是“子列没拆干净”或者“把极限点和整个极限混了”。直接在页面里选中那一步继续问我就行。
